\documentclass[11pt]{article}

\usepackage[english]{babel}
\usepackage[letterpaper,top=2cm,bottom=2cm,left=3cm,right=3cm]{geometry}
\usepackage{amsmath,amssymb,amsthm,mathtools}
\usepackage{booktabs}
\usepackage{array}
\usepackage{enumitem}
\usepackage[numbers,sort&compress]{natbib}
\usepackage{xcolor}
\definecolor{ethblue}{rgb}{0.129, 0.361, 0.686}
\usepackage[colorlinks=true,linkcolor=ethblue,citecolor=ethblue,urlcolor=ethblue]{hyperref}
\usepackage[noabbrev]{cleveref}

% ---------------------------------------------------------------------------
% Theorem environments
% ---------------------------------------------------------------------------
\theoremstyle{plain}
\newtheorem{theorem}{Theorem}[section]
\newtheorem{proposition}[theorem]{Proposition}
\newtheorem{corollary}[theorem]{Corollary}
\newtheorem{lemma}[theorem]{Lemma}
\theoremstyle{definition}
\newtheorem{definition}[theorem]{Definition}
\newtheorem{assumption}[theorem]{Assumption}
\newtheorem{dataAssumption}[theorem]{Data Assumption}
\newtheorem{example}[theorem]{Example}
\theoremstyle{remark}
\newtheorem{remark}[theorem]{Remark}

\crefname{dataAssumption}{Data Assumption}{Data Assumptions}
\crefname{assumption}{Assumption}{Assumptions}

% ---------------------------------------------------------------------------
% Notation
% ---------------------------------------------------------------------------
\newcommand{\RR}{\mathbb{R}}
\newcommand{\NN}{\mathbb{N}}
\newcommand{\E}{\mathbb{E}}
\newcommand{\Prob}{\mathcal{P}}
\newcommand{\UU}{\mathcal{U}}
\newcommand{\VV}{\mathcal{V}}
\newcommand{\XX}{\mathcal{X}}
\newcommand{\KK}{\mathcal{K}}
\newcommand{\Dd}{\mathrm{d}}
\newcommand{\define}{\coloneqq}
\newcommand{\muP}{\mu_{\mathsf{P}}}
\newcommand{\KL}[2]{D_{\mathrm{KL}}\!\left(#1\,\Vert\,#2\right)}
\newcommand{\DE}{D_{\mathrm{E}}}
\newcommand{\IP}[2]{\left\langle #1,\, #2 \right\rangle}
\newcommand{\abs}[1]{\left\lvert #1 \right\rvert}
\newcommand{\norm}[1]{\left\lVert #1 \right\rVert}
\newcommand{\indicator}{\mathbf{1}}
\newcommand{\sdot}{s_{\mathrm{dot}}}
\newcommand{\sncos}{s_{\mathrm{ncos}}}
\newcommand{\sltwo}{s_{L^2}}
\newcommand{\ESS}{\mathrm{ESS}}
\newcommand{\Jvu}{J_{v\vert u}}
\newcommand{\Juv}{J_{u\vert v}}
\newcommand{\Jtwo}{J_{2}}
\newcommand{\Jjoint}{J_{\mathrm{joint}}}

\title{Admissibility and Tail Behaviour in Contrastive Forecasting}
\author{Aleksandr Beliaev\\ \small University of Toronto \quad \texttt{aleksandr.beliaev@mail.utoronto.ca}}
\date{October 2026}

\begin{document}
\maketitle

\begin{abstract}
 Contrastive learning can be read as a methodology for learning a conditional
 distribution: a learned score tilts the marginal distribution of one modality,
 and the tilted measure is the model. Adopting this population-level viewpoint,
 developed by Baptista, Stuart and Tran, we ask which design choices in a
 contrastive forecaster survive in the limit of infinite data and flexible
 encoders: the loss (one-way, two-way, or joint), the statistical divergence
 that it minimizes, or the score. We first show that every loss variant, under
 every divergence, has the same population-level target, the log density ratio
 up to a conditioning-dependent constant; combined with the universal
 approximation theorems of Stuart and Wolf, the loss and the divergence drop
 out for flexible encoders. We then show that the score does not drop out, and
 that it matters precisely through its tail behaviour. A score that is bounded
 above defines a valid model for every data distribution, and the forecast
 obtained by reweighting a finite bank of observed futures converges to that
 model as the bank grows. A score that grows without bound in the future,
 such as the unnormalized inner product, admits no valid model on heavy-tailed
 data, even with ReLU encoders, and its bank forecast collapses onto a single
 future as the bank grows. This identifies the tail condition in existing
 non-compact universal approximation results as necessary rather than
 technical. Numerical experiments on synthetic and real time series confirm
 both halves and add one refinement: in practice a ceiling on the score is
 necessary but not sufficient, and the score must also not reward extreme
 futures.
\end{abstract}

\noindent\textbf{Keywords:} Contrastive Learning, CLIP, Probabilistic
Forecasting, Importance Sampling, Heavy Tails, Model Misspecification

% ===========================================================================
\section{Introduction}
\label{sec:intro}
% ===========================================================================

Contrastive learning trains two encoders, one for each of two paired data
modalities, so that matched pairs receive a high score and mismatched pairs a
low one. In the forecasting setting considered here, the two modalities are
the past $u$ and the future $v$ of a time series. A recent line of work
\cite{baptista2025mathematical,stuart2026expressivity} views such algorithms
probabilistically: the learned score $s(u,v)$ defines a probability model for
$v$ given $u$, obtained by reweighting the marginal distribution of futures by
$e^{s}$, and the contrastive losses used in practice are statistical
divergences between this model and the data. This paper studies which design
choices in the resulting forecaster matter, and why.

In Subsection~\ref{ssec:BALR} we provide background and review related work.
Subsection~\ref{ssec:CAO} states our contributions and describes the
organization of the paper. Subsection~\ref{ssec:NOT} fixes notation.

\subsection{Background and Literature Review}
\label{ssec:BALR}

\paragraph{The probabilistic framework.}
Our starting point is the framework of Baptista, Stuart and Tran
\cite{baptista2025mathematical}. There, a score defines a tilting of the
product of the marginals, the CLIP loss and its one-way and joint variants are
identified as Kullback--Leibler (KL) divergences between the data and the
tilted model \cite[Thm.~3.2, Thm.~3.8, Rem.~3.12]{baptista2025mathematical},
generalizations to other divergences are formulated
\cite[Opt.~Probs.~3.4 and 3.6, Rem.~3.5]{baptista2025mathematical}, and a
forecast is obtained by reweighting a finite bank of observed futures
\cite[Eq.~(31)]{baptista2025mathematical}. For Gaussian data and linear
encoders every minimizer is available in closed form, and the choice of loss
and of score visibly changes the resulting forecast
\cite[Sec.~5]{baptista2025mathematical}.

\paragraph{Expressivity.}
Stuart and Wolf \cite{stuart2026expressivity} adopt the same population-level
view and study expressivity. They show that the joint KL divergence between
the data and the model bounds the conditional losses from above
\cite[Thms.~2.5, 2.6]{stuart2026expressivity}, that uniform approximation of
the log density ratio lifts to approximation of the measure in KL
\cite[Thms.~3.2, 3.3]{stuart2026expressivity}, and that two-tower CLIP, with
unnormalized or normalized encoders, is a universal approximator in KL on
compact domains \cite[Thms.~4.1, 4.2]{stuart2026expressivity}. Their
extension to non-compact domains truncates the approximating score outside a
compact set; the underlying estimate requires the approximant to be bounded
by a constant $B$ outside that set and carries an error term of order $e^{B}$
\cite[Prop.~A.3, Thm.~3.3]{stuart2026expressivity}. They work with KL
throughout and list other divergences, such as those built from the energy
score, as future work.

\paragraph{The optimal score.}
That the optimal InfoNCE critic is the pointwise mutual information plus a
function of the conditioning variable is standard
\cite{oord2018cpc,poole2019variational}. For the symmetric CLIP loss, Uesaka
et al.\ \cite{uesaka2025weighted} show that the optimal similarity is the
pointwise mutual information and argue that inner products of
finite-dimensional embeddings may fail to reach it. This is the closest
relative of the misspecification view taken in \Cref{sec:invariance}, but it
concerns a single loss rather than agreement or disagreement between
variants.

\paragraph{Reweighting candidates, and its failure.}
Sordoni et al.\ \cite{sordoni2021decomposed} note that resampling candidates
with softmax-normalized critic scores approximates the conditional
distribution, exactly as the pool of candidates grows; this is the bank
forecast studied here. Weight collapse of self-normalized importance sampling
is classical: in high dimension it is studied by Bengtsson, Bickel and Li
\cite{bengtsson2008curse}, heavy-tailed weights are diagnosed by
Pareto-smoothed importance sampling \cite{vehtari2024pareto}, and the
required sample size is of order $e^{D_{\mathrm{KL}}}$ by Chatterjee and
Diaconis \cite{chatterjee2018sample}. That an exponential tilt
$e^{\langle a, v\rangle}$ requires a finite moment generating function is
textbook. Our results in \Cref{sec:tails} apply these facts to contrastive
scores; we claim no new probability.

\paragraph{Bounded critics and normalization.}
In mutual-information estimation, Song and Ermon \cite{song2020limitations}
clip the critic to tame the variance of partition-function estimates. In
representation learning, normalized embeddings are standard because they
remove the scale degree of freedom that lets the softmax become arbitrarily
sharp \cite{chen2020simclr,wang2020understanding}, and negative squared
Euclidean distance logits are used in practice and match or exceed cosine
CLIP \cite{chou2024embedding}. To our knowledge none of these works ties
boundedness to the \emph{existence} of the model, or to the behaviour of a
reweighted bank.

\paragraph{Analog forecasting.}
Forecasting by weighting historical futures is the analog ensemble of
weather forecasting \cite{dellemonache2013analog} and kernel analog
forecasting \cite{zhao2016analog}. A contrastive bank forecast is an analog
forecast whose kernel, $e^{s(u,v)}$, is learned.

\subsection{Contribution and Paper Organization}
\label{ssec:CAO}

The primary contributions of this paper are a population-level account of
which design choices in a contrastive forecaster matter, together with
numerical evidence on synthetic and real time series. In detail, we make the
following contributions, structured to simultaneously describe the paper
organization:
\begin{enumerate}[start=1,label={(C\arabic*)}]
 \item \emph{Loss and divergence invariance} (\Cref{sec:invariance}).
       We show that the one-way, two-way and joint losses, under any
       divergence, vanish exactly when the score equals the log density ratio
       up to, respectively, a function of the past, a function of the future,
       or a constant. Consequently, all of them agree at the population level
       whenever the score family is well-specified. Combining this with the
       results of Stuart and Wolf, every KL-based loss variant is
       simultaneously driven to zero by flexible two-tower encoders on
       compact domains.
 \item \emph{Loss differences as misspecification} (\Cref{ssec:gaussian}).
       In the Gaussian, linear setting of \cite{baptista2025mathematical} we
       show that each score family omits a different ingredient of the log
       density ratio, which explains the predictions of
       \cite{baptista2025mathematical} as artefacts of misspecification. We
       prove that linear $L^2$ scores cannot reach zero two-way or joint loss
       on any correlated Gaussian, and give the scalar optimum in closed form.
 \item \emph{Admissibility is governed by the tails of the score}
       (\Cref{ssec:admissibility}).
       A score bounded above yields a valid model for every data
       distribution. A score linear in the future yields no valid model on
       heavy-tailed data unless it ignores the past, and ReLU encoders do not
       remove this obstruction. Boundedness has a price: a bounded score can
       be exact only when the density ratio is bounded.
 \item \emph{Convergence and collapse of the bank forecast}
       (\Cref{ssec:bank}).
       We show that the bank forecast converges as the bank grows if and only
       if the score is admissible, that it collapses onto a single future
       when the scores over the bank are regularly varying, and that a
       two-sided bound guarantees a uniform effective sample size at the cost
       of capped sharpness. This shows that the tail condition in
       \cite[Prop.~A.3]{stuart2026expressivity} is necessary rather than a
       convenience of the proof.
 \item \emph{Numerical experiments} (\Cref{sec:experiments}).
       On synthetic Gaussian, heavy-tailed, warped and bimodal processes, and
       on real weather data, we confirm the loss invariance for flexible
       encoders, the failure of unbounded scores on heavy tails, and the
       predicted turnaround of the unbounded forecast as the bank grows. The
       experiments add one refinement: a ceiling on the score is necessary but
       not sufficient, and the score must also not reward extreme futures.
\end{enumerate}
\Cref{sec:conclusion} concludes and lists limitations and open questions.
Proofs are collected in \Cref{app:proofs}.

\subsection{Notation}
\label{ssec:NOT}

For a measurable space we write $\Prob(\cdot)$ for the set of probability
measures on it. The past $u$ takes values in $\UU\subseteq\RR^{d_u}$, the
future $v$ in $\VV\subseteq\RR^{d_v}$, and $\XX\define\UU\times\VV$. The data
distribution is $\mu\in\Prob(\XX)$, with marginals $\mu_u$ and $\mu_v$ and
product of marginals $\muP\define\mu_u\otimes\mu_v$. Conditionals are written
$\mu_{v\vert u}(\cdot\vert u)$ and $\mu_{u\vert v}(\cdot\vert v)$. Euclidean
inner products and norms are $\IP{\cdot}{\cdot}$ and $\abs{\cdot}$. For
symmetric matrices, $A\preceq B$ denotes the Loewner order. $\delta_x$ is the
Dirac measure at $x$, and $\indicator_A$ the indicator of a set $A$.

% ===========================================================================
\section{Probabilistic Formulation}
\label{sec:setting}
% ===========================================================================

In this section we specify the data model (\Cref{ssec:problem}), the model
defined by a score (\Cref{ssec:model}), the loss functions
(\Cref{ssec:losses}), and the finite-sample forecast that is used in
practice (\Cref{ssec:bankdef}). The presentation follows
\cite{baptista2025mathematical,stuart2026expressivity}.

\subsection{Problem Formulation}
\label{ssec:problem}

\begin{assumption}[Data Distribution]
 \label{ass:data}
 The data generating measure $\mu\in\Prob(\XX)$ satisfies the following.
 \begin{enumerate}[start=1,label={(A\arabic*)}]
  \item \label{ass:marginals}\emph{Marginal distributions.}
        The marginals $\mu_u\in\Prob(\UU)$ and $\mu_v\in\Prob(\VV)$ are
        absolutely continuous with respect to Lebesgue measure.
  \item \label{ass:joint}\emph{Joint distribution.}
        The measures $\mu$ and $\muP$ are mutually absolutely continuous, with
        Radon--Nikodym derivative
        \begin{equation}
         r(u,v)\define\frac{\Dd\mu}{\Dd\muP}(u,v)
         =\frac{p(v\vert u)}{p(v)}.
         \label{eq:ratio}
        \end{equation}
  \item \label{ass:MI}\emph{Finite mutual information.}
        $\log r\in L^1(\XX,\mu)$, so that
        $I(u;v)\define\E_{\mu}\log r<\infty$.
 \end{enumerate}
\end{assumption}

The function $\log r$ is the pointwise mutual information: it measures how
much more likely the future $v$ is after observing the past $u$ than it is a
priori. It is the tilting function of \cite{stuart2026expressivity}, and it
is the object that every loss studied below aims to recover.

\subsection{The Model Defined by a Score}
\label{ssec:model}

\begin{definition}[Score and Tilted Model]
 \label{def:model}
 A \emph{score} is a measurable function $s:\XX\to\RR$. It defines the
 joint and conditional measures
 \begin{equation}
  \nu^{s}(\Dd u,\Dd v)\define\frac{e^{s(u,v)}}{Z}\,\muP(\Dd u,\Dd v),
  \qquad
  \nu^{s}_{v\vert u}(\Dd v\vert u)\define\frac{e^{s(u,v)}}{Z(u)}\,\mu_v(\Dd v),
  \label{eq:model}
 \end{equation}
 with normalizing constants
 \begin{equation}
  Z\define\int_{\XX}e^{s}\,\Dd\muP,\qquad
  Z(u)\define\int_{\VV}e^{s(u,v)}\,\mu_v(\Dd v),\qquad
  Z'(v)\define\int_{\UU}e^{s(u,v)}\,\mu_u(\Dd u),
  \label{eq:normalizers}
 \end{equation}
 and $\nu^{s}_{u\vert v}$ defined analogously with $Z'(v)$
 \cite[Eqs.~(13)--(16)]{baptista2025mathematical}. When $Z<\infty$, the
 conditional of $\nu^{s}$ given $u$ is $\nu^{s}_{v\vert u}$.
\end{definition}

In words, the forecast for a past $u$ takes the distribution of all futures
and reweights each one by $e^{\text{score}}$. The construction requires the
normalizing constants to be finite, which motivates the following
definition; it is central to \Cref{sec:tails}.

\begin{definition}[Admissible Score]
 \label{def:admissible}
 A score $s$ is \emph{admissible for the forecasting direction} if
 $Z(u)<\infty$ for $\mu_u$-almost every $u$, and \emph{admissible} if, in
 addition, $Z'(v)<\infty$ for $\mu_v$-almost every $v$ and $Z<\infty$. If
 a normalizing constant is infinite, the corresponding model does not exist,
 and we assign the value $+\infty$ to every loss that involves it.
\end{definition}

\begin{example}[Scores]
 \label{ex:scores}
 Let $f:\UU\to\RR^{d_z}$ and $g:\VV\to\RR^{d_z}$ be encoders of the past and
 the future, and let $\kappa>0$ be a learnable scale (the inverse
 temperature, $\kappa=1/\tau$ in \cite{stuart2026expressivity}). We consider
 \begin{equation}
  \sdot=\kappa\IP{f(u)}{g(v)},\qquad
  \sncos=\kappa\IP{\frac{f(u)}{\abs{f(u)}}}{\frac{g(v)}{\abs{g(v)}}},\qquad
  \sltwo=-\frac{\kappa}{2}\abs{f(u)-g(v)}^{2}.
  \label{eq:scores}
 \end{equation}
 The unnormalized inner product $\sdot$ is the score of
 \cite[Thm.~4.1]{stuart2026expressivity}, the normalized cosine $\sncos$ is
 that of CLIP \cite{radford2021clip} and of
 \cite[Thm.~4.2]{stuart2026expressivity}, and $\sltwo$ is the $L^2$ score of
 \cite{baptista2025mathematical}. Note that $\sltwo\le0$ and
 $\abs{\sncos}\le\kappa$, whereas $\sdot$ is unbounded whenever $g$ is.
\end{example}

\subsection{Population-Level Loss Functions}
\label{ssec:losses}

\begin{definition}[Divergence]
 \label{def:divergence}
 A \emph{divergence} is a map $D:\Prob\times\Prob\to[0,\infty]$ with
 $D(p,q)=0$ if and only if $p=q$.
\end{definition}

The KL divergence is a divergence. So is the energy distance, on
distributions with a finite first moment,
\begin{equation}
 \DE(p,q)\define2\,\E\abs{X-Y}-\E\abs{X-X'}-\E\abs{Y-Y'},
 \qquad X,X'\sim p,\;Y,Y'\sim q\text{ independent},
 \label{eq:energy}
\end{equation}
as is the maximum mean discrepancy with a characteristic kernel. The energy
distance is also our evaluation metric: the expected energy score
\cite{gneiting2007strictly} of a forecast $q$, minus that of the truth $p$,
equals $\tfrac12\DE(p,q)$.

\begin{definition}[Contrastive Losses]
 \label{def:losses}
 For a divergence $D$, weights $\lambda_u,\lambda_v>0$ and a score $s$,
 define the \emph{one-way}, \emph{reverse one-way}, \emph{two-way} and
 \emph{joint} losses
 \begin{equation}
  \begin{aligned}
   \Jvu^{D}(s)    & \define\E_{u\sim\mu_u}\,D\bigl(\mu_{v\vert u}(\cdot\vert u),\nu^{s}_{v\vert u}(\cdot\vert u)\bigr), &
   \Juv^{D}(s)    & \define\E_{v\sim\mu_v}\,D\bigl(\mu_{u\vert v}(\cdot\vert v),\nu^{s}_{u\vert v}(\cdot\vert v)\bigr),   \\
   \Jtwo^{D}(s)   & \define\lambda_u\Jvu^{D}(s)+\lambda_v\Juv^{D}(s),                                                    &
   \Jjoint^{D}(s) & \define D\bigl(\mu,\nu^{s}\bigr).
  \end{aligned}
  \label{eq:losses}
 \end{equation}
 When $D=D_{\mathrm{KL}}$ we drop the superscript.
\end{definition}

\begin{remark}[Relation to CLIP and InfoNCE]
 \label{rem:clip}
 With $D=D_{\mathrm{KL}}$, expanding \eqref{eq:losses} with \eqref{eq:model}
 gives $\Jvu(s)=I(u;v)-\E_{\mu}[s]+\E_{\mu_u}[\log Z(u)]$. Replacing the
 measures by empirical measures yields the InfoNCE loss for $\Jvu$, the CLIP
 loss for $\Jtwo$ with $\lambda_u=\lambda_v=\tfrac12$, and a single softmax
 over all pairs for $\Jjoint$
 \cite[Thm.~3.2, Thm.~3.8, Rem.~3.12]{baptista2025mathematical}; see also
 \cite[Sec.~2.2]{stuart2026expressivity}. For general $D$ the losses
 correspond to \cite[Opt.~Probs.~3.4 and 3.6]{baptista2025mathematical}; they
 can be trained from samples alone when $D$ is the KL divergence, an energy
 distance, or a maximum mean discrepancy
 \cite[Rem.~3.5]{baptista2025mathematical}.
\end{remark}

\subsection{The Bank Forecast}
\label{ssec:bankdef}

The model \eqref{eq:model} involves the full marginal $\mu_v$. In practice it
is replaced by a finite collection of observed futures.

\begin{dataAssumption}
 \label{dass:bank}
 We have access to a \emph{bank} $v_1,\ldots,v_N\overset{\mathrm{iid}}{\sim}\mu_v$
 of futures, independent of the data used to train the score.
\end{dataAssumption}

Following \cite[Eq.~(31)]{baptista2025mathematical}, the \emph{bank forecast}
for a past $u$ is
\begin{equation}
 \nu^{N}_{v\vert u}(\cdot\vert u)\define\sum_{j=1}^{N}w_j\,\delta_{v_j},\qquad
 w_j\define\frac{e^{s(u,v_j)}}{\sum_{k=1}^{N}e^{s(u,v_k)}},\qquad
 \ESS\define\Bigl(\sum_{j=1}^{N}w_j^{2}\Bigr)^{-1}.
 \label{eq:bank}
\end{equation}
This is self-normalized importance sampling with proposal $\mu_v$ and target
$\nu^{s}_{v\vert u}$. The effective sample size $\ESS\in[1,N]$ counts how
many bank futures the forecast effectively uses; $\ESS\approx1$ means that
the forecast has collapsed onto a single future.

% ===========================================================================
\section{Loss and Divergence Invariance}
\label{sec:invariance}
% ===========================================================================

In this section we show that every loss variant, under every divergence,
aims at the same object, the log density ratio. The variants differ only in
which part of it they leave free, and the divergences differ only in how they
measure a miss. Differences between them are therefore a symptom of a score
family that cannot represent $\log r$.

\begin{proposition}[Common Target]
 \label{prop:target}
 Let $\mu$ satisfy \Cref{ass:data}, let $D$ be any divergence, and let $s$ be
 an admissible score. Then
 \begin{enumerate}[label=(\roman*)]
  \item $\Jvu^{D}(s)=0$ if and only if $s(u,v)=\log r(u,v)+c(u)$ for some
        function $c$;
  \item $\Juv^{D}(s)=0$ if and only if $s(u,v)=\log r(u,v)+c'(v)$ for some
        function $c'$;
  \item $\Jtwo^{D}(s)=0$ if and only if $\Jjoint^{D}(s)=0$ if and only if
        $s=\log r+\mathrm{const}$.
 \end{enumerate}
 All identities hold $\muP$-almost everywhere. The right-hand sides do not
 involve $D$.
\end{proposition}
\begin{proof}
 See \Cref{app:target}.
\end{proof}

\begin{definition}[Well-Specified Family]
 \label{def:wellspecified}
 A family $\mathcal{S}$ of scores is \emph{well-specified} for a loss if that
 loss attains the value $0$ on $\mathcal{S}$. By \Cref{prop:target}, this means
 that $\mathcal{S}$ contains some $\log r+c(u)$ for the one-way loss, and
 $\log r+\mathrm{const}$ for the two-way and joint losses, whatever the
 divergence.
\end{definition}

\begin{corollary}[Agreement]
 \label{cor:agreement}
 If $\mathcal{S}$ contains $\log r+\mathrm{const}$, then for every divergence
 $D$ the one-way, two-way (for any $\lambda_u,\lambda_v>0$) and joint losses
 all have minimizers whose forecast is exactly the true conditional,
 $\nu^{s}_{v\vert u}=\mu_{v\vert u}$. Neither the loss variant nor the
 divergence can be distinguished at the population level.
\end{corollary}

The converse fails: \Cref{ssec:gaussian} contains a family that cannot
represent $\log r$, yet for which the one-way and two-way losses agree.

\subsection{Flexible Encoders}
\label{ssec:flexible}

\Cref{cor:agreement} requires the family to contain $\log r$ exactly. For
flexible encoders, the results of Stuart and Wolf provide the approximate
analogue for the KL divergence.

\begin{theorem}[Approximate Agreement, KL]
 \label{thm:approxagreement}
 Let $\mu$ satisfy \Cref{ass:data} with $\UU$ and $\VV$ compact, and let
 the encoders be drawn from a parametric family satisfying the
 $C^0$-universal approximation assumption of
 \cite[Sec.~3]{stuart2026expressivity}. Then for every $\epsilon>0$
 there exist a latent dimension $d_z$ and encoders $f,g$ such that
 $s=\sdot$ satisfies
 \begin{equation*}
  \Jjoint(s)<\epsilon,\qquad\Jvu(s)<\epsilon,\qquad\Juv(s)<\epsilon,
  \qquad\Jtwo(s)<(\lambda_u+\lambda_v)\,\epsilon.
 \end{equation*}
 The same holds for $s=\sncos$, with $\kappa$ chosen large enough.
\end{theorem}
\begin{proof}
 See \Cref{app:approxagreement}.
\end{proof}

A single flexible score therefore makes every KL-based loss variant small
simultaneously. This is the first half of our account: once the encoders
are flexible, the loss variant drops out. The statement is restricted to
compact domains, on which continuous encoders produce bounded scores;
\Cref{sec:tails} studies what happens when that boundedness is lost.

\Cref{thm:approxagreement} does not transfer to other divergences, because
closeness in KL does not imply closeness in other divergences. The simplest
case is a forecast that is correct except for a small amount of mass moved far
away.

\begin{proposition}[KL Is Blind to How Far Misplaced Mass Goes]
 \label{prop:klblind}
 Let $p$ be a distribution on $\RR^{d}$ with a finite mean and no atom at
 $a\in\RR^d$, and let $q_\varepsilon=(1-\varepsilon)p+\varepsilon\delta_a$.
 Then
 \begin{equation*}
  \KL{p}{q_\varepsilon}\le-\log(1-\varepsilon)\quad\text{for every }a,
 \end{equation*}
 whereas
 \begin{equation*}
  \tfrac12\DE(p,q_\varepsilon)
  =\tfrac12\varepsilon^{2}\bigl(2\,\E\abs{X-a}-\E\abs{X-X'}\bigr)
  \xrightarrow[\abs{a}\to\infty]{}\infty.
 \end{equation*}
\end{proposition}
\begin{proof}
 See \Cref{app:klblind}.
\end{proof}

Moving a fraction $\varepsilon$ of the forecast to an arbitrarily extreme
future costs at most about $\varepsilon$ in KL, while its cost in energy
score grows linearly with the distance. A model trained with KL therefore has
almost no incentive to keep mass out of the far tails, whereas an evaluation
by energy score penalizes exactly that. The divergences thus agree exactly at
the optimum (\Cref{cor:agreement}) and approximately in the bulk
(\Cref{thm:approxagreement}), but an approximation that is good in one may be
poor in another in the tails. Training with the energy distance itself, as
\cite[Rem.~3.5]{baptista2025mathematical} allows, would penalize tail
over-weighting directly.

\subsection{Rigid Families: the Gaussian Case}
\label{ssec:gaussian}

When the family is not flexible, the loss variants come apart. The Gaussian,
linear setting of \cite[Sec.~5]{baptista2025mathematical} makes this
explicit. There, every minimizer is available in closed form, and four
predictions are relevant here:
\begin{enumerate}[label=(P\arabic*),ref=P\arabic*]
 \item \label{P1} the inner-product score gets the mean of the forecast right
       but its spread wrong, which remains at the unconditional level
       \cite[Cor.~5.2]{baptista2025mathematical};
 \item \label{P2} the one-way $L^2$ score gets both mean and spread right
       \cite[Thm.~5.3, Cor.~5.4]{baptista2025mathematical};
 \item \label{P3} for the inner-product score, the one-way and two-way losses
       give the same answer \cite[Thm.~5.1]{baptista2025mathematical};
 \item \label{P4} the joint loss shrinks the forecast towards the mean and is
       a worse forecaster than the one-way loss
       \cite[Thm.~5.6, Cor.~5.7]{baptista2025mathematical}.
\end{enumerate}
We now show that each prediction is a consequence of misspecification.

Let $\mu$ be a centred Gaussian with covariance blocks $C_{uu},C_{uv},C_{vv}$,
joint covariance $C$, joint precision $P=C^{-1}$, and
$C_0\define\mathrm{blockdiag}(C_{uu},C_{vv})$, the covariance of $\muP$. Then
\begin{equation}
 \log r(u,v)=-\tfrac12u^{\top}(P_{uu}-C_{uu}^{-1})u\;-\;u^{\top}P_{uv}v
 \;-\;\tfrac12v^{\top}(P_{vv}-C_{vv}^{-1})v+\mathrm{const}.
 \label{eq:gaussianlogr}
\end{equation}
There are three ingredients: a past--future interaction, which sets the mean
of the forecast; a future-only term, which sets its spread; and a past-only
term. With linear encoders every score is of the form
$s=-\tfrac12u^{\top}Ru+u^{\top}Av-\tfrac12v^{\top}Sv$, and the families
differ in which ingredients they can set (\Cref{tab:families}).

\begin{table}[ht]
 \centering
 \small
 \begin{tabular}{lll}
  \toprule
  family                      & allowed $(R,A,S)$                                                          & well-specified for          \\
  \midrule
  inner product               & $R=S=0$                                                                    & nothing (if $C_{uv}\ne0$)   \\
  $L^2$                       & $\begin{pmatrix}R&-A\\-A^{\top}&S\end{pmatrix}=[G,-H]^{\top}[G,-H]\succeq0$ & one-way loss only           \\
  inner product + quadratics  & $R,S$ symmetric, $A$ free                                                  & every loss                  \\
  \bottomrule
 \end{tabular}
 \caption{Linear score families on Gaussian data. For $L^2$, $f(u)=Gu$ and
  $g(v)=Hv$, with $\kappa$ absorbed.}
 \label{tab:families}
\end{table}

\Cref{tab:families} recovers the predictions. For \ref{P1}: the inner
product has no future-only term, so the covariance of its forecast is
$C_{vv}$ whatever $A$ is. For \ref{P4}: the joint loss also needs the
past-only term, which the inner product lacks; moreover, the joint normalizer
is finite only if $\norm{C_{uu}^{1/2}AC_{vv}^{1/2}}<1$, which is the origin of
the shrinkage function of \cite{baptista2025mathematical}. For \ref{P2}: the
one-way $L^2$ family is well-specified, because $S\succeq0$ can match the
future-only term. The two-way and joint $L^2$ losses are not.

\begin{proposition}[Linear $L^2$ Is Too Rigid for the Two-Way and Joint Losses]
 \label{prop:l2rigid}
 Under linear $L^2$ scores, every model joint covariance satisfies
 $(C_0^{-1}+Q)^{-1}\preceq C_0$, with $Q=[G,-H]^{\top}[G,-H]\succeq0$. The
 data satisfy $C\not\preceq C_0$ whenever $C_{uv}\neq0$. Hence the two-way
 and joint losses cannot reach zero.
\end{proposition}
\begin{proof}
 See \Cref{app:l2rigid}.
\end{proof}

In words, an $L^2$ model can couple past and future only by shrinking both
marginals, while real data couple them without shrinking.

\begin{proposition}[Scalar Closed Form]
 \label{prop:scalar}
 For unit-variance scalars with correlation $\rho\in(0,1)$, the joint-$L^2$
 optimum forecasts
 $\mathcal{N}\bigl(\tfrac{\rho}{2-\rho}u,\tfrac{2(1-\rho)}{2-\rho}\bigr)$,
 against the truth $\mathcal{N}(\rho u,1-\rho^{2})$. The mean is shrunk and the
 forecast is over-confident. The fraction of the mutual information $I(u;v)$
 that it loses tends to $3/8$ as $\rho\to0$, and to $0$ as $\rho\to1$.
\end{proposition}
\begin{proof}
 See \Cref{app:scalar}.
\end{proof}

\begin{table}[ht]
 \centering
 \small
 \setlength{\tabcolsep}{4pt}
 \begin{tabular}{lc|ccc|ccc|c}
  \toprule
                           &          & \multicolumn{3}{c|}{inner product} & \multicolumn{3}{c|}{$L^2$} & + quad. \\
  setting                  & $I(u;v)$ & 1-way & 2-way & joint              & 1-way & 2-way & joint      & all three            \\
  \midrule
  $\rho=0.1$               & 0.005    & 0.005 & 0.005 & 0.005              & 0     & 0.294 & 0.333      & 0                    \\
  $\rho=0.3$               & 0.047    & 0.046 & 0.046 & 0.051              & 0     & 0.219 & 0.251      & 0                    \\
  $\rho=0.5$               & 0.144    & 0.131 & 0.131 & 0.157              & 0     & 0.148 & 0.170      & 0                    \\
  $\rho=0.7$               & 0.337    & 0.272 & 0.272 & 0.323              & 0     & 0.081 & 0.092      & 0                    \\
  $\rho=0.9$               & 0.830    & 0.512 & 0.512 & 0.571              & 0     & 0.021 & 0.023      & 0                    \\
  GP of \Cref{sec:experiments} & 1.627 & 0.704 & 0.704 & 0.746 & 0  & 0.003 & 0.003      & 0                    \\
  \bottomrule
 \end{tabular}
 \caption{Fraction of the mutual information $I(u;v)$ (in nats) lost by the
  forecast at the population optimum of each KL loss: $0$ is perfect, $1$ is
  the unconditional forecast. Values are exact, computed from Gaussian KL
  divergences without training or sampling. The last row uses the exact
  covariances of the synthetic Gaussian process of
  \Cref{sec:experiments}.}
 \label{tab:infoloss}
\end{table}

\Cref{tab:infoloss} makes the picture quantitative, and two points stand out.
First, each rigid family fails in its own regime. The error of the inner
product lies in the spread, which matters most when dependence is strong; the
error of two-way and joint $L^2$ lies in the coupling-by-shrinking
constraint, which matters most when dependence is weak. Second, adding the
missing ingredients removes every difference: giving the inner product free
past-only and future-only quadratic terms makes all three losses exact (last
column). This is the flexibility argument of \Cref{ssec:flexible} in
miniature.

% ===========================================================================
\section{Admissibility and the Tails of the Score}
\label{sec:tails}
% ===========================================================================

Flexibility concerns representing $\log r$ where the data are. It says
nothing about how the score behaves far from the data, and that is where the
choice of score continues to matter. The non-compact universal approximation
result of \cite[Thm.~3.3]{stuart2026expressivity} rests on an estimate
\cite[Prop.~A.3]{stuart2026expressivity} which requires the approximating
score to satisfy $\abs{s}\le B$ outside a compact set $\KK$, and whose error
contains a term of order $e^{B}$; the theorem then enforces the bound by
truncating the approximant to zero outside $\KK$. For a two-tower score and
$\KK=\KK_u\times\KK_v$, truncation corresponds to encoders that vanish
outside compact sets, since
$\indicator_{\KK}\IP{f}{g}=\IP{\indicator_{\KK_u}f}{\indicator_{\KK_v}g}$.
Standard neural encoders do not vanish. In this section we show that the tail
condition is not a convenience of the proof: without it the model need not
exist, and the bank forecast collapses.

\subsection{Admissibility}
\label{ssec:admissibility}

\begin{proposition}[Bounded Above Implies Admissible]
 \label{prop:bounded}
 If $s\le c$, then $Z(u)\le e^{c}$, $Z'(v)\le e^{c}$ and $Z\le e^{c}$ for
 every data distribution $\mu$. In particular, $\sltwo$ (with $c=0$) and
 $\sncos$ (with $c=\kappa$) are admissible for every $\mu$, with any
 encoders.
\end{proposition}
\begin{proof}
 See \Cref{app:bounded}.
\end{proof}

\begin{proposition}[A Score Linear in the Future Needs Light Tails]
 \label{prop:linear}
 Let $s(u,v)=\IP{a(u)}{Hv}+b(u)$. Then $Z(u)=e^{b(u)}M_{\mu_v}(H^{\top}a(u))$,
 where $M_{\mu_v}$ is the moment generating function of $\mu_v$. If $\mu_v$
 is a multivariate Student-$t$ distribution, with any number of degrees of
 freedom, then $M_{\mu_v}(w)=\infty$ for every $w\neq0$. Hence the linear
 inner-product score is admissible for the forecasting direction only if
 $H^{\top}a(u)=0$ for $\mu_u$-almost every $u$, that is, only if it ignores
 the past.
\end{proposition}
\begin{proof}
 See \Cref{app:linear}.
\end{proof}

\begin{proposition}[ReLU Encoders Do Not Restore Admissibility]
 \label{prop:relu}
 Let $s=\sdot$ with $g$ a ReLU network, and let $g_\infty$ denote the same
 network with its biases removed. If $\IP{f(u)}{g_\infty(v_0)}>0$ for some
 direction $v_0$, and $\mu_v$ is an elliptical Student-$t$ distribution, then
 $Z(u)=\infty$.
\end{proposition}
\begin{proof}
 See \Cref{app:relu}.
\end{proof}

The inner-product score is therefore admissible at $u$ only if it does not
grow in any direction of the future, and nothing in training enforces this.
A bounded output layer on $g$, such as $\tanh$, makes the score admissible by
\Cref{prop:bounded} without normalizing anything; \Cref{ssec:banksweep}
shows that this is not sufficient in practice.

Boundedness is not free. Combining \Cref{prop:target} with an upper bound on
the score constrains the density ratios that the family can represent
exactly.

\begin{proposition}[Bounded Scores Represent Only Bounded Ratios]
 \label{prop:price}
 Let $s\le c$ be admissible and let $D$ be any divergence.
 \begin{enumerate}[label=(\roman*)]
  \item If $\Jvu^{D}(s)=0$, then for $\mu_u$-almost every $u$,
        $r(u,\cdot)\le e^{c}/Z(u)$ holds $\mu_v$-almost everywhere.
  \item If $\Jjoint^{D}(s)=0$, then $r\le e^{c}/Z$ holds $\muP$-almost
        everywhere.
  \item If $\abs{s}\le\kappa$ and $\Jjoint^{D}(s)=0$, then
        $e^{-2\kappa}\le r\le e^{2\kappa}$ holds $\muP$-almost everywhere.
 \end{enumerate}
\end{proposition}
\begin{proof}
 See \Cref{app:price}.
\end{proof}

For a correlated Gaussian, $r$ is unbounded on $\XX$, so no score bounded
above, however flexible, is well-specified for the joint or two-way loss;
\Cref{prop:l2rigid} is a special case. This does not contradict
\Cref{thm:approxagreement}: the approximants constructed in
\cite[Thm.~3.3]{stuart2026expressivity} are themselves bounded, being
continuous on $\KK$ and zero outside it. Boundedness thus costs exactness but
not approximability in KL. For the normalized cosine, part (iii) is the
counterpart of the observation in \cite[Thm.~4.2]{stuart2026expressivity}
that universality requires the temperature $1/\kappa$ to tend to zero.

\subsection{Convergence and Collapse of the Bank Forecast}
\label{ssec:bank}

Admissibility is not a technicality: it is exactly the condition under which
the bank forecast \eqref{eq:bank} works.

\begin{theorem}[Admissibility Is Equivalent to Bank Convergence]
 \label{thm:bank}
 Fix $u$ and let the bank satisfy \Cref{dass:bank}. If $Z(u)<\infty$, then
 $\nu^{N}_{v\vert u}\to\nu^{s}_{v\vert u}$ weakly, almost surely, as
 $N\to\infty$, and
 $\ESS/N\to1/\E_{\mu_v}\bigl[(\Dd\nu^{s}_{v\vert u}/\Dd\mu_v)^{2}\bigr]$ when
 this moment is finite. If $Z(u)=\infty$, then the empirical normalizer
 $\frac1N\sum_{j}e^{s(u,v_j)}$ diverges and there is no limit to converge
 to.
\end{theorem}
\begin{proof}
 See \Cref{app:bank}.
\end{proof}

\begin{proposition}[Collapse]
 \label{prop:collapse}
 If the scores $X_j=s(u,v_j)$ of the bank futures have a regularly varying
 right tail, $\mathbb{P}(X>x)=x^{-\alpha}L(x)$ with $L$ slowly varying, then
 $\max_j w_j\to1$ in probability as $N\to\infty$. This covers the linear
 inner-product score on Student-$t$ futures, with $\alpha$ equal to the
 degrees of freedom.
\end{proposition}
\begin{proof}
 See \Cref{app:collapse}.
\end{proof}

\begin{proposition}[Bounded on Both Sides: Guaranteed Bank, Capped Sharpness]
 \label{prop:twosided}
 If $\abs{s}\le\kappa$, then
 $e^{-2\kappa}\le\Dd\nu^{s}_{v\vert u}/\Dd\mu_v\le e^{2\kappa}$. Hence
 $\KL{\nu^{s}_{v\vert u}}{\mu_v}\le2\kappa$ and, in the limit $N\to\infty$,
 $\ESS/N\ge e^{-2\kappa}$, uniformly in $u$. The price is that no forecast can
 be more than $e^{2\kappa}$ times as concentrated as the marginal
 $\mu_v$, or carry more than $2\kappa$ nats of information.
\end{proposition}
\begin{proof}
 See \Cref{app:twosided}.
\end{proof}

The two outcomes are therefore precise. With a bounded score the bank
forecast converges to the model as the bank grows. With an unbounded score on
heavy-tailed data it collapses onto the single most extreme future, and a
larger bank only supplies more extreme futures to collapse onto.
\Cref{tab:fixedscore} checks this for fixed scores, with no training, and
\Cref{tab:hierarchy} summarizes the hierarchy.

\begin{table}[ht]
 \centering
 \small
 \begin{tabular}{rccc}
  \toprule
  bank size $N$ & linear, Gaussian & linear, Student-$t$ ($\nu=3$) & $L^2$, Student-$t$ ($\nu=3$) \\
  \midrule
  $10^{2}$      & 0.081\ (0.41)          & 0.240\ (0.30)                       & 0.018\ (0.79)                      \\
  $10^{3}$      & 0.016\ (0.38)          & 0.455\ (0.038)                      & 0.002\ (0.79)                      \\
  $10^{4}$      & 0.003\ (0.37)          & 0.760\ (0.0007)                     & $<10^{-3}$\ (0.79)                 \\
  $10^{5}$      & 0.001\ (0.37)          & 0.902\ ($10^{-5}$)                  & $<10^{-4}$\ (0.79)                 \\
  $10^{6}$      & $<10^{-4}$\ (0.37)     & 0.954\ ($10^{-6}$)                  & $<10^{-5}$\ (0.79)                 \\
  \bottomrule
 \end{tabular}
 \caption{Largest bank weight (mean over 200 banks), with $\ESS/N$ in
  brackets, for a fixed score $s=v$ or $s=-\tfrac12(v-1)^{2}$ on
  unit-variance futures. For the Gaussian, \Cref{thm:bank} gives
  $\ESS/N\to e^{-1}\approx0.368$.}
 \label{tab:fixedscore}
\end{table}

\begin{table}[ht]
 \centering
 \begin{tabular}{>{\raggedright\arraybackslash}p{3.0cm}>{\raggedright\arraybackslash}p{3.3cm}>{\raggedright\arraybackslash}p{3.5cm}>{\raggedright\arraybackslash}p{3.3cm}}
  \toprule
                         & unbounded ($\sdot$)                 & bounded above ($\sltwo$)                                               & bounded on both sides ($\sncos$) \\
  \midrule
  model exists           & only with light-tailed futures      & always                                                                 & always                           \\
  bank forecast converges& only with light-tailed futures      & always                                                                 & always                           \\
  bank size needed       & can be infinite (collapse)          & finite, $\approx e^{\KL{\nu_{v\vert u}}{\mu_v}}$ \cite{chatterjee2018sample} & uniformly $\lesssim e^{2\kappa}$ \\
  sharpness              & unlimited                           & unlimited                                                              & capped by $\kappa$               \\
  \bottomrule
 \end{tabular}
 \caption{The boundedness hierarchy.}
 \label{tab:hierarchy}
\end{table}

\begin{remark}[Why Training Does Not Avoid Inadmissibility]
 \label{rem:training}
 At the population level an inadmissible parameter has infinite loss, so an
 exact minimizer would avoid it. Training, however, uses the empirical loss,
 which is always finite: a batch of $256$ futures almost never contains the
 extreme futures that make $Z(u)$ infinite. Training can therefore settle on
 an informative parameter whose score grows in the future, and the problem
 appears only when the forecast meets a large bank. For the linear inner
 product on Student-$t$ data, \Cref{prop:linear} shows that every informative
 parameter is of this kind. \Cref{prop:klblind} gives a second reason: even
 when the model exists, KL training barely charges for misplaced tail mass.
\end{remark}

% ===========================================================================
\section{Numerical Experiments}
\label{sec:experiments}
% ===========================================================================

We test the predictions of \Cref{sec:invariance,sec:tails} on synthetic
processes, each of which relaxes one assumption of the Gaussian setting, and
on real weather data. \Cref{ssec:setup} describes the setup.
\Cref{ssec:lossvariants} tests the loss invariance of
\Cref{sec:invariance}, \Cref{ssec:heavytails,ssec:banksweep} test the tail
results of \Cref{sec:tails}, and \Cref{ssec:finite} discusses finite data.

\subsection{Setup}
\label{ssec:setup}

\paragraph{Evaluation.}
Every contrastive model is evaluated only through its own score, by building
the bank forecast \eqref{eq:bank} from a bank of real futures not used in
training. Forecasts are scored with the energy score \cite{gneiting2007strictly},
a proper scoring rule for multivariate probabilistic forecasts that rewards
both accuracy and calibrated uncertainty. On synthetic data each domain has
an exact optimal forecaster (oracle), and we report the energy score as a
percentage above that of the oracle ($0\%$ is perfect). On real data we report
\emph{skill}, the fractional reduction in energy score against a reference
forecast (higher is better).

\paragraph{Synthetic domains.}
Past and future are windows of $50$ points of a smooth Gaussian process
observed with noise. Each synthetic domain has a parameter whose zero point
is this Gaussian process, the setting of \cite{baptista2025mathematical}:
\begin{itemize}
 \item \emph{heavy tails}: an elliptical Student-$t$ process with $\nu$
       degrees of freedom ($\nu=\infty$ is the Gaussian process; $\nu=3$ is
       the heaviest used);
 \item \emph{bent GP}: a monotone warp of the Gaussian process, giving
       skewed futures and a nonlinear relation between past and future;
 \item \emph{double-well}: a stochastic process that jumps between two
       levels, so that the future can have two separate likely outcomes.
\end{itemize}
Each synthetic setting uses $200$k training pairs, a bank of $50$k held-out
futures, $500$ test cases and $2$ seeds, so that differences below about
$1\%$ are noise.

\paragraph{Real data.}
We use one year (2020) of 10-minute readings from the Jena weather station. We
forecast the next 6\,h of temperature from the last day of six variable
groups, and the next 6\,h of rain from past pressure, wind and moisture.

\paragraph{Encoders.}
Linear encoders reproduce the setting of \cite{baptista2025mathematical}.
Flexible encoders are two-layer MLPs of width $256$ with GELU activations and
a linear shortcut.

\subsection{Loss Variants}
\label{ssec:lossvariants}

\paragraph{Linear encoders: the loss and the score both matter.}
On the Gaussian process, the inner-product score is $+5.6\%$ above perfect
whichever direction it is trained in (\ref{P3}), and $+9.8\%$ when trained
jointly (\ref{P4}). Its $90\%$ bands contain the truth $98\%$ of the time on
the hardest third of cases, i.e.\ they are too wide (\ref{P1}). The $L^2$
score is at $+0.0\%$ in all three variants, consistent with the last row of
\Cref{tab:infoloss}. The linear normalized cosine is poor everywhere
($+12$--$23\%$), because it sees the past only through its direction and
loses its level.

\paragraph{Flexible encoders: the loss variant mostly stops mattering.}
With MLP encoders the range between the variants collapses
(\Cref{tab:variants}), as \Cref{thm:approxagreement} predicts, and \ref{P1}
fades: the coverage of the inner product on the hardest cases drops from
$98\%$ to $89\%$. The best models (joint $L^2$, normalized cosine) are within
about $1\%$ of perfect on every domain. The differences that remain are
small and consistently favour the joint loss; we return to them in
\Cref{ssec:finite}.

\begin{table}[ht]
 \centering
 \begin{tabular}{llcccc}
  \toprule
  domain                       & encoder & inner prod.\ 1-way & 2-way & joint & range \\
  \midrule
  Gaussian process             & linear  & $+5.6$  & $+5.6$  & $+9.8$  & 4.2  \\
                               & MLP     & $+2.7$  & $+1.4$  & $+0.3$  & 2.4  \\
  bent GP (full warp)          & linear  & $+37.4$ & $+40.0$ & $+25.2$ & 14.8 \\
                               & MLP     & $+5.9$  & $+3.9$  & $+0.6$  & 5.3  \\
  double-well ($\lambda=1.5$)  & linear  & $+1.3$  & $+2.6$  & $+2.7$  & 1.4  \\
                               & MLP     & $+0.6$  & $+0.3$  & $+0.5$  & 0.3  \\
  \bottomrule
 \end{tabular}
 \caption{Energy score above the oracle (\%) for the inner-product score
  trained with each loss variant. The range across variants shrinks with
  flexible encoders.}
 \label{tab:variants}
\end{table}

\subsection{Unbounded Scores on Heavy Tails}
\label{ssec:heavytails}

Every catastrophic failure we observed is an unbounded score on heavy tails
(\Cref{tab:heavy}). The linear inner product places essentially all its
weight on single bank futures (minimum $\ESS$ of $1.0$). MLP encoders do not
rescue it, in line with \Cref{prop:relu}: MLP inner products are still $+7$
to $+23\%$ above perfect and still collapse onto single futures for some
inputs. On rain, which is mostly zero with occasional bursts, the
unnormalized inner product scored $1.4$--$2.4$ times worse than climatology,
while $L^2$ and normalized cosine retained the shape ``mostly dry, sometimes a
burst''. Normalization helps exactly where tails are heavy, and only there.

\begin{table}[ht]
 \centering
 \begin{tabular}{lccc}
  \toprule
  heavy tails, $\nu=3$ & inner product, 2-way             & normalized cosine & joint $L^2$ \\
  \midrule
  linear               & $\sim10\times$ worse than oracle & $+23.0$           & $+4.5$      \\
  MLP                  & $+23.3$                          & $+0.9$            & $+0.7$      \\
  \bottomrule
 \end{tabular}
 \caption{Energy score above the oracle (\%) on the heavy-tailed domain.}
 \label{tab:heavy}
\end{table}

\subsection{Bank-Size Sweep}
\label{ssec:banksweep}

The two experiments in this subsection test \Cref{thm:bank} and \Cref{prop:collapse}
directly. Each model is trained as in \Cref{ssec:setup} ($15$k steps, batch
size $256$) and then frozen; only the bank changes. The forecast is rebuilt
from held-out banks of size $N=100$ up to $N=200{,}000$, with up to $3$
disjoint banks per size, on the same $500$ test pasts. We record the energy
score above the oracle, the largest bank weight and, for collapsed forecasts
(largest weight $>0.5$), how extreme the winning future is, measured by its
percentile by size within the bank. We use $2$ seeds on the heavy-tailed
domain ($\nu=3$) and $1$ seed on the Gaussian control.

\paragraph{Does a larger bank help or hurt?}
\Cref{thm:bank} and \Cref{prop:collapse} predict that bounded scores improve and level
off as the bank grows, and that the unbounded inner product turns around and
gets worse. The prediction holds (\Cref{tab:banksweep}). On heavy tails, the
unbounded inner product is best at a small bank and then gets steadily worse:
with MLP encoders from $+7\%$ at $3$k futures to $+46\%$ at $200$k, and with
linear encoders from $+33\%$ to over $+3000\%$. Every bounded score improves
and levels off. On the Gaussian control the same unbounded score also
improves and levels off (at $+5.7\%$ for linear encoders, matching
\Cref{ssec:lossvariants}), so the turnaround is caused by the tails and not
by the score alone. The collapse diagnostic shows the mechanism directly:
with linear encoders, $20\%$ of the forecasts of the inner product at
$N=200$k put more than half their weight on one future, and that future is
the most extreme in the bank ($99.998$th percentile by size, seed 1).

\paragraph{Is a ceiling enough?}
\Cref{prop:bounded} implies that a $\tanh$ output layer on the future encoder
makes the inner product admissible without normalizing anything. The result
is only a partial fix (\Cref{tab:banksweep}, second row of each block). It is
far better than the uncapped inner product ($+103\%$ against $+3358\%$ with
linear encoders, $+7.5\%$ against $+46\%$ with MLPs, at $N=200$k), but it
still turns around: in one MLP seed it improves to $+1.2\%$, in the other it
rises from $+8\%$ to $+14\%$, and with linear encoders it rises from $+22\%$
to $+103\%$. Its forecasts rarely collapse onto one future, but the futures
they favour are still the most extreme in the bank ($99.99$th percentile,
linear, seed 1).

This does not contradict \Cref{sec:tails}. The ceiling makes the model
exist, so the bank forecast converges; but it converges to a model that
over-weights extremes. A $\tanh$ score rises \emph{towards} its ceiling as
the future becomes more extreme, so extreme futures in the favoured
direction receive the highest score. A small bank contains few of them and
hides the problem; a large bank exposes it.

\begin{table}[ht]
 \centering
 \footnotesize
 \setlength{\tabcolsep}{3.5pt}
 \begin{tabular}{llrrrrrr}
  \toprule
                            & model                         & $N=100$          & 1k               & 3k              & 30k             & 100k    & 200k    \\
  \midrule
  heavy tails, MLP          & inner product (unbounded)     & $+33.6$          & $+8.7$           & $\mathbf{+7.1}$ & $+17.5$         & $+34.5$ & $+46.3$ \\
                            & inner product, $\tanh$ future & $+37.0$          & $+9.7$           & $+6.7$          & $\mathbf{+3.2}$ & $+5.4$  & $+7.5$  \\
                            & $L^2$ one-way                 & $+36.8$          & $+11.8$          & $+8.0$          & $+4.5$          & $+3.8$  & $+2.9$  \\
                            & $L^2$ joint                   & $+31.5$          & $+6.2$           & $+3.2$          & $+0.9$          & $+0.7$  & $+0.6$  \\
                            & normalized cosine             & $+41.8$          & $+8.0$           & $+3.5$          & $+1.0$          & $+0.8$  & $+0.8$  \\
  \midrule
  heavy tails, linear       & inner product (unbounded)     & $\mathbf{+33.2}$ & $+38.7$          & $+70.2$         & $+1253$         & $+2705$ & $+3358$ \\
                            & inner product, $\tanh$ future & $+28.0$          & $\mathbf{+22.1}$ & $+24.5$         & $+64.1$         & $+92.3$ & $+102.6$\\
                            & $L^2$ one-way                 & $+22.8$          & $+8.1$           & $+5.8$          & $+4.6$          & $+4.3$  & $+4.2$  \\
                            & $L^2$ joint                   & $+21.5$          & $+7.8$           & $+5.7$          & $+4.5$          & $+4.5$  & $+4.5$  \\
                            & normalized cosine             & $+25.4$          & $+23.5$          & $+23.2$         & $+23.1$         & $+23.0$ & $+23.1$ \\
  \midrule
  Gaussian, MLP             & inner product (unbounded)     & $+13.7$          & $+2.9$           & $+1.5$          & $+2.0$          & $+1.8$  & $+1.8$  \\
  Gaussian, linear          & inner product (unbounded)     & $+9.9$           & $+6.4$           & $+6.1$          & $+5.8$          & $+5.7$  & $+5.7$  \\
  \bottomrule
 \end{tabular}
 \caption{Energy score above the oracle (\%) as the bank grows, for frozen
  models. Mean of $2$ seeds ($1$ seed for the Gaussian control). Bold marks the best
  bank size where a curve turns around.}
 \label{tab:banksweep}
\end{table}

\paragraph{What the two experiments show together.}
The scores that work either \emph{penalize} extreme futures ($\sltwo$
decreases as $g(v)$ moves away from $f(u)$) or are \emph{indifferent} to how
extreme a future is ($\sncos$ depends only on the direction of $g(v)$). It is
not the height of the ceiling that matters: the learned scale of the
normalized cosine is about $\kappa\approx72$, so the guarantee of
\Cref{prop:twosided} is vacuous, and yet it is the best model here. The
underlying reason is \Cref{rem:training}: training cannot correct the tails.
A batch of $256$ futures almost never contains one from the far tail, so how
the score treats extreme futures is determined by the architecture, not
learned from data. A ceiling is necessary, since without one the model does
not exist; in practice the score must also not \emph{reward} extremity. We
have not stated or proved this refinement as a theorem.

\subsection{Finite Data: the Loss Matters Again}
\label{ssec:finite}

The results of \Cref{sec:invariance} are statements about infinite data.
With finite data the loss variant can matter again, through over-fitting
rather than through what it aims at. On weather data, MLPs over-fit within a
few hundred to a few thousand steps. There, the joint $L^2$ loss clearly beats
the one-way $L^2$ loss (\Cref{tab:weather}), although one-way $L^2$ is the
well-specified variant.

\begin{table}[ht]
 \centering
 \begin{tabular}{lccc}
  \toprule
  skill ($2$ seeds)                      & $L^2$ 1-way & $L^2$ 2-way & $L^2$ joint \\
  \midrule
  rain, linear (vs climatology)          & $-0.106$    & $+0.011$    & $+0.012$    \\
  temperature, MLP (vs persistence)      & $+0.152$    & $+0.243$    & $+0.280$    \\
  temperature, linear (vs persistence)   & $+0.352$    & $+0.344$    & $+0.383$    \\
  \bottomrule
 \end{tabular}
 \caption{Skill of the $L^2$ score on Jena weather data (higher is better).}
 \label{tab:weather}
\end{table}

The population analysis bounds what joint and two-way $L^2$ can lose to
rigidity: in the scalar Gaussian, at most about $3/8$ (joint) or $1/3$
(two-way) of $I(u;v)$ (\Cref{prop:scalar}). On rain, $I(u;v)$ is tiny (the
best model beats climatology by $1\%$), so that cost is negligible, and the
observed gaps must come from finite data. One-way $L^2$ has nothing that
prevents it from becoming over-sharp, and on rain its forecasts collapsed
(minimum $\ESS$ of $1.6$ with linear encoders, median $\ESS$ of $3$ with
MLPs). Two-way and joint $L^2$ do have such a check. By a Schur complement,
$S\succeq A^{\top}R^{+}A$, so a sharp forecast must come with a sharp
backward conditional, which the two-way loss penalizes; the joint loss also
scores the model marginals, which a collapsed forecast distorts. On this
reading, the small bias of the joint loss buys a regularizer against
over-sharpening. This predicts that the deficit of the one-way loss shrinks
as the training data grow, as it does on the plentiful synthetic data of
\Cref{ssec:lossvariants}. It is a hypothesis, not a result.

% ===========================================================================
\section{Conclusion and Future Work}
\label{sec:conclusion}
% ===========================================================================

We have viewed a contrastive forecaster as a learned tilting of the marginal
distribution of futures, and asked which of its design choices survive in the
population limit. The loss and the divergence do not: every variant aims at
the log density ratio, and for flexible encoders on compact domains a single
score drives every KL-based variant to zero. The score does survive, and
precisely through its tails. A score bounded above defines a valid model for
every data distribution, and its bank forecast converges; a score that grows
in the future defines no model on heavy-tailed data, even with ReLU encoders,
and its bank forecast collapses onto the most extreme future. This identifies
the tail bound in the non-compact universal approximation result of
\cite{stuart2026expressivity} as necessary rather than technical.
Experiments on synthetic and real time series confirm both halves, and
suggest that a bounded score must, in addition, not reward extreme futures.

\paragraph{Limitations.}
\Cref{sec:invariance,sec:tails} are population-level statements (infinite
data, infinitely many negatives); the finite-sample statements are
\Cref{thm:bank} and \Cref{prop:collapse}, the experiments of \Cref{ssec:banksweep}, and
the hypothesis of \Cref{ssec:finite}. The information losses in
\Cref{tab:infoloss} are KL divergences, while the experiments use the energy
score, which is less sensitive to errors in spread; the two should agree in
sign but not in size. \Cref{prop:relu} is proved for ReLU networks, while our
MLPs use GELU, which grows like ReLU at infinity; we expect the same
conclusion but have not written out the proof. Rain is physically bounded,
so its true $Z(u)$ is finite and \Cref{prop:collapse} does not apply
literally; what applies is its finite-$N$ version, a small $\ESS/N$ whenever
the scores over the bank are heavy-tailed. All experiments use $2$ seeds,
except the Gaussian control of \Cref{ssec:banksweep} ($1$ seed), and the
collapse diagnostic was recorded for seed 1 only.

\paragraph{Open questions.}
First, the divergence part of \Cref{sec:invariance} is proved but not
tested: every model here was trained with KL. The natural test is to train
with the energy-distance loss permitted by
\cite[Rem.~3.5]{baptista2025mathematical}; by \Cref{prop:klblind} we expect
this to matter most for scores that reward extreme futures. Second, the
``does not reward extremity'' refinement of \Cref{ssec:banksweep} is an
empirical reading of two experiments on one domain and should be stated and
proved as a theorem. Third, it remains to establish non-compact universal
approximation for an architecture whose score is bounded by construction,
such as the $L^2$ or normalized-cosine two-tower families, rather than by
truncation. Finally, the weak-dependence prediction of \Cref{prop:scalar} at
finite data, the quadratic-augmented inner product (\Cref{tab:infoloss},
last column) in training, and the bank sweep on rain remain untested.

% ===========================================================================
\appendix
\section{Proofs}
\label{app:proofs}
% ===========================================================================

\subsection{Proof of \texorpdfstring{\Cref{prop:target}}{Proposition 3.1}}
\label{app:target}
Because $D(p,q)\ge0$ with equality only if $p=q$, each loss vanishes exactly
when the distributions that it compares are equal. The remainder of the
argument is therefore about equality of distributions and does not involve
$D$.

(i) $\Jvu^{D}(s)=0$ if and only if $\nu^{s}_{v\vert u}(\cdot\vert u)=\mu_{v\vert u}(\cdot\vert u)$
for $\mu_u$-almost every $u$. Comparing densities with respect to $\mu_v$,
this reads $e^{s(u,\cdot)}/Z(u)=r(u,\cdot)$, that is,
$s=\log r+\log Z(u)$. Conversely, $s=\log r+c(u)$ gives
$Z(u)=e^{c(u)}\int r(u,v)\,\mu_v(\Dd v)=e^{c(u)}$, and hence the true
conditional. (ii) is symmetric. (iii) Since $\lambda_u,\lambda_v>0$,
$\Jtwo^{D}=0$ if and only if (i) and (ii) both hold, so that $c(u)=c'(v)$ for
$\muP$-almost every $(u,v)$. A function of $u$ alone that equals a function
of $v$ alone almost everywhere under a product measure is almost everywhere
constant. For the joint loss, $D(\mu,\nu^{s})=0$ if and only if
$\nu^{s}=\mu$, if and only if $e^{s}/Z=r$, that is, $s=\log r+\log Z$.
\qed

\subsection{Proof of \texorpdfstring{\Cref{thm:approxagreement}}{Theorem 3.4}}
\label{app:approxagreement}
By \cite[Thm.~4.1]{stuart2026expressivity} there exist $d_z$ and continuous
encoders $f,g$ such that $s=\IP{f}{g}$ (with $\kappa=1$ absorbed into the
encoders) satisfies $\KL{\mu}{\nu^{s}}<\epsilon$. Since $\UU,\VV$ are compact
and $s$ is continuous, $s$ is bounded, so $s$ is admissible and $\mu\ll\nu^{s}$.
The measure $\nu^{s}$ has density $e^{s}/Z$ with respect to $\muP$, so its
$u$-marginal has density $Z(u)/Z$ with respect to $\mu_u$, and its
conditional given $u$ is $\nu^{s}_{v\vert u}$ of \eqref{eq:model}; likewise
for $v$. Applying \cite[Thm.~2.6]{stuart2026expressivity} with $m=2$ to the
pairs $(v,u)$ and $(u,v)$ gives $\Jvu(s)\le\KL{\mu}{\nu^{s}}$ and
$\Juv(s)\le\KL{\mu}{\nu^{s}}$, and the bound on $\Jtwo$ follows. For
$\sncos$, use \cite[Thm.~4.2]{stuart2026expressivity} in place of
\cite[Thm.~4.1]{stuart2026expressivity}, with $\kappa=1/\tau$.
\qed

\subsection{Proof of \texorpdfstring{\Cref{prop:klblind}}{Proposition 3.5}}
\label{app:klblind}
Since $p$ has no atom at $a$, $q_\varepsilon\ge(1-\varepsilon)p$ on the
support of $p$, so $\Dd p/\Dd q_\varepsilon\le1/(1-\varepsilon)$ $p$-almost
everywhere and $\KL{p}{q_\varepsilon}=\int\log(\Dd p/\Dd q_\varepsilon)\,\Dd p\le-\log(1-\varepsilon)$.
For the energy distance, write $m=\E\abs{X-X'}$ and $A=\E\abs{X-a}$. With
$Y\sim q_\varepsilon$, $\E\abs{X-Y}=(1-\varepsilon)m+\varepsilon A$ and
$\E\abs{Y-Y'}=(1-\varepsilon)^{2}m+2\varepsilon(1-\varepsilon)A$.
Substituting, $\DE=2(1-\varepsilon)m+2\varepsilon A-m-(1-\varepsilon)^{2}m-2\varepsilon(1-\varepsilon)A=\varepsilon^{2}(2A-m)$,
and $A\ge\abs{a}-\E\abs{X}\to\infty$.
\qed

\subsection{Proof of \texorpdfstring{\Cref{prop:l2rigid}}{Proposition 3.6}}
\label{app:l2rigid}
The model joint has density proportional to $e^{s}\,\mathcal{N}(0,C_0)$ with
$s=-\tfrac12z^{\top}Qz$ and $z=(u,v)$, that is, it is
$\mathcal{N}\bigl(0,(C_0^{-1}+Q)^{-1}\bigr)$. Since
$C_0^{-1}+Q\succeq C_0^{-1}$, inverting reverses the order, so
$(C_0^{-1}+Q)^{-1}\preceq C_0$. On the data side,
$C_0-C=\begin{pmatrix}0&-C_{uv}\\-C_{uv}^{\top}&0\end{pmatrix}$ has
eigenvalues $\pm\sigma_i(C_{uv})$, so $C\preceq C_0$ forces $C_{uv}=0$. Zero
two-way or joint loss requires the model joint to equal $\mu$
(\Cref{prop:target}(iii)), which is impossible when $C_{uv}\neq0$. The
inequality $S\succeq A^{\top}R^{+}A$ used in \Cref{ssec:finite} is the
generalized Schur complement criterion for
$\begin{pmatrix}R&-A\\-A^{\top}&S\end{pmatrix}\succeq0$.
\qed

\subsection{Proof of \texorpdfstring{\Cref{prop:scalar}}{Proposition 3.7}}
\label{app:scalar}
With $s=-\tfrac12(gu-hv)^{2}$ and unit variances,
$Q=\begin{pmatrix}g^{2}&-gh\\-gh&h^{2}\end{pmatrix}$ and, up to constants,
$2\Jjoint=g^{2}+h^{2}-2gh\rho-\log(1+g^{2}+h^{2})$. For fixed $t=g^{2}+h^{2}$
the cross term is minimized at $g=h$, leaving $t(1-\rho)-\log(1+t)$, which is
minimized at $t=\rho/(1-\rho)$. The forecast then has precision
$1+t/2=(2-\rho)/(2(1-\rho))$ and mean coefficient $\rho/(2-\rho)$. Expanding
for small $\rho$ gives $\Jvu=\tfrac{3}{16}\rho^{2}+O(\rho^{3})$ against
$I(u;v)=\tfrac12\rho^{2}+O(\rho^{4})$, so the ratio tends to $3/8$. At
$\rho=0.5$ the closed form gives $0.1699$, matching \Cref{tab:infoloss}.
\qed

\subsection{Proof of \texorpdfstring{\Cref{prop:bounded}}{Proposition 4.1}}
\label{app:bounded}
$Z(u)=\int e^{s(u,v)}\,\mu_v(\Dd v)\le e^{c}$, and likewise for $Z'(v)$ and
$Z$. For $L^2$, $-\tfrac{\kappa}{2}\abs{f-g}^{2}\le0$; for the normalized
cosine, $\abs{\IP{\bar f}{\bar g}}\le1$ by the Cauchy--Schwarz inequality.
\qed

\subsection{Proof of \texorpdfstring{\Cref{prop:linear}}{Proposition 4.2}}
\label{app:linear}
$Z(u)=e^{b(u)}\,\E\exp\IP{H^{\top}a(u)}{v}$ by definition. A multivariate
$t$ distribution is elliptical, so for $w\neq0$ the projection $\IP{w}{v}$ is
a scaled univariate $t$ variable, whose density decays polynomially; hence
$\E e^{\IP{w}{v}}=\infty$.
\qed

\subsection{Proof of \texorpdfstring{\Cref{prop:relu}}{Proposition 4.3}}
\label{app:relu}
Write $g=W_L\sigma(\cdots\sigma(W_1v+b_1)\cdots)+b_L$ with
$\sigma=\mathrm{ReLU}$. The ReLU is $1$-Lipschitz and positively homogeneous,
so by induction over the layers
$\abs{g(v)-g_\infty(v)}\le\sum_{\ell}\norm{b_\ell}\prod_{k>\ell}\norm{W_k}\eqqcolon B<\infty$,
and $g_\infty(tv)=t\,g_\infty(v)$ for $t\ge0$. Let
$\delta=\IP{f(u)}{g_\infty(v_0/\abs{v_0})}>0$. By continuity there is an open
cone $\KK\ni v_0$ on which $\IP{f(u)}{g_\infty(v/\abs{v})}\ge\delta'>0$, so
on $\KK$, $s(u,v)\ge\kappa\bigl(\delta'\abs{v}-\abs{f(u)}B\bigr)$. An
elliptical $t$ distribution gives every open cone a polynomial tail,
$\mu_v(\KK\cap\{\abs{v}>t\})\asymp t^{-\nu}$, so
$\int_{\KK}e^{\kappa\delta'\abs{v}}\,\Dd\mu_v=\infty$.
\qed

\subsection{Proof of \texorpdfstring{\Cref{prop:price}}{Proposition 4.4}}
\label{app:price}
(i) By \Cref{prop:target}(i) and its proof, $r(u,\cdot)=e^{s(u,\cdot)}/Z(u)\le e^{c}/Z(u)$.
(ii) By \Cref{prop:target}(iii), $r=e^{s}/Z\le e^{c}/Z$. (iii) If
$\abs{s}\le\kappa$ then $Z\in[e^{-\kappa},e^{\kappa}]$, so
$r=e^{s}/Z\in[e^{-2\kappa},e^{2\kappa}]$.
\qed

\subsection{Proof of \texorpdfstring{\Cref{thm:bank}}{Theorem 4.5}}
\label{app:bank}
For bounded continuous $\varphi$, the strong law of large numbers gives
$\frac1N\sum_{j}e^{s(u,v_j)}\varphi(v_j)\to\int e^{s}\varphi\,\Dd\mu_v$ and
$\frac1N\sum_{j}e^{s(u,v_j)}\to Z(u)$ almost surely, since both integrands
are integrable when $Z(u)<\infty$. The ratio converges to
$\int\varphi\,\Dd\nu^{s}_{v\vert u}$, and a countable convergence-determining
class gives weak convergence. With $\rho=e^{s}/Z(u)$,
$N\sum_{j}w_j^{2}=\frac{\frac1N\sum_j\rho(v_j)^{2}}{(\frac1N\sum_j\rho(v_j))^{2}}\to\E_{\mu_v}[\rho^{2}]$.
If $Z(u)=\infty$, the strong law for nonnegative random variables gives
$\frac1N\sum_{j}e^{s(u,v_j)}\to\infty$.
\qed

\subsection{Proof of \texorpdfstring{\Cref{prop:collapse}}{Proposition 4.6}}
\label{app:collapse}
Let $X_{(1)}\ge X_{(2)}$ be the two largest scores and $b_N$ the
$1-1/N$ quantile of $X$, so that $b_N=N^{1/\alpha}\ell(N)$ with $\ell$ slowly
varying. By the extreme-value limit for regularly varying tails,
$(X_{(1)},X_{(2)})/b_N$ converges in distribution to
$(\Gamma_1^{-1/\alpha},\Gamma_2^{-1/\alpha})$, the first two points of a
Poisson process, which are almost surely distinct. So for every
$\varepsilon>0$ there is $c>0$ with
$\mathbb{P}(X_{(1)}-X_{(2)}\ge cb_N)\ge1-\varepsilon$ for large $N$, and on
that event
$1-\max_j w_j\le\sum_{j\neq(1)}e^{X_j-X_{(1)}}\le Ne^{-cb_N}\to0$, because
$b_N$ grows polynomially.
\qed

\subsection{Proof of \texorpdfstring{\Cref{prop:twosided}}{Proposition 4.7}}
\label{app:twosided}
$\abs{s}\le\kappa$ gives $Z(u)\in[e^{-\kappa},e^{\kappa}]$, so
$\Dd\nu^{s}_{v\vert u}/\Dd\mu_v=e^{s}/Z(u)\in[e^{-2\kappa},e^{2\kappa}]$.
Then
$\KL{\nu^{s}_{v\vert u}}{\mu_v}=\E_{\nu}\log(\Dd\nu/\Dd\mu_v)\le2\kappa$ and
$\E_{\mu_v}[(\Dd\nu/\Dd\mu_v)^{2}]=\E_{\nu}[\Dd\nu/\Dd\mu_v]\le e^{2\kappa}$;
the limit of $\ESS/N$ follows from \Cref{thm:bank}.
\qed

\paragraph{Reproducibility.}
\Cref{tab:infoloss,tab:fixedscore} are produced by
\begin{center}\texttt{test\_bed/theory/misspecification.py},\end{center}
and the experiments of \Cref{ssec:banksweep} by
\begin{center}\texttt{test\_bed/analysis/bank\_sweep.py}.\end{center}

% ===========================================================================
\begin{thebibliography}{99}

\bibitem{baptista2025mathematical}
R.~Baptista, A.~M. Stuart, and S.~Tran.
\newblock A mathematical perspective on contrastive learning.
\newblock \emph{arXiv preprint arXiv:2505.24134}, 2025.

\bibitem{stuart2026expressivity}
A.~M. Stuart and F.~Wolf.
\newblock Expressivity in multimodal contrastive learning.
\newblock \emph{arXiv preprint arXiv:2608.17203}, 2026.

\bibitem{radford2021clip}
A.~Radford, J.~W. Kim, C.~Hallacy, A.~Ramesh, G.~Goh, S.~Agarwal, G.~Sastry,
 A.~Askell, P.~Mishkin, J.~Clark, G.~Krueger, and I.~Sutskever.
\newblock Learning transferable visual models from natural language
 supervision.
\newblock In \emph{International Conference on Machine Learning}, 2021.

\bibitem{oord2018cpc}
A.~van~den Oord, Y.~Li, and O.~Vinyals.
\newblock Representation learning with contrastive predictive coding.
\newblock \emph{arXiv preprint arXiv:1807.03748}, 2018.

\bibitem{poole2019variational}
B.~Poole, S.~Ozair, A.~van~den Oord, A.~Alemi, and G.~Tucker.
\newblock On variational bounds of mutual information.
\newblock In \emph{International Conference on Machine Learning}, 2019.

\bibitem{uesaka2025weighted}
T.~Uesaka et~al.
\newblock Weighted point set embedding for multimodal contrastive learning
 toward optimal similarity metric.
\newblock In \emph{International Conference on Learning Representations},
 2025. arXiv:2404.19228.

\bibitem{sordoni2021decomposed}
A.~Sordoni et~al.
\newblock Decomposed mutual information estimation for contrastive
 representation learning.
\newblock In \emph{International Conference on Machine Learning}, 2021.
 arXiv:2106.13401.

\bibitem{song2020limitations}
J.~Song and S.~Ermon.
\newblock Understanding the limitations of variational mutual information
 estimators.
\newblock In \emph{International Conference on Learning Representations},
 2020.

\bibitem{chen2020simclr}
T.~Chen, S.~Kornblith, M.~Norouzi, and G.~Hinton.
\newblock A simple framework for contrastive learning of visual
 representations.
\newblock In \emph{International Conference on Machine Learning}, 2020.

\bibitem{wang2020understanding}
T.~Wang and P.~Isola.
\newblock Understanding contrastive representation learning through alignment
 and uniformity on the hypersphere.
\newblock In \emph{International Conference on Machine Learning}, 2020.

\bibitem{chou2024embedding}
J.~C.-C. Chou et~al.
\newblock Embedding geometries of contrastive language-image pre-training.
\newblock In \emph{ECCV Workshops}, 2024. arXiv:2409.13079.

\bibitem{bengtsson2008curse}
T.~Bengtsson, P.~Bickel, and B.~Li.
\newblock Curse-of-dimensionality revisited: collapse of the particle filter
 in very large scale systems.
\newblock \emph{IMS Collections}, 2:316--334, 2008.

\bibitem{vehtari2024pareto}
A.~Vehtari, D.~Simpson, A.~Gelman, Y.~Yao, and J.~Gabry.
\newblock Pareto smoothed importance sampling.
\newblock \emph{Journal of Machine Learning Research}, 25, 2024.
 arXiv:1507.02646.

\bibitem{chatterjee2018sample}
S.~Chatterjee and P.~Diaconis.
\newblock The sample size required in importance sampling.
\newblock \emph{The Annals of Applied Probability}, 28(2):1099--1135, 2018.

\bibitem{dellemonache2013analog}
L.~Delle~Monache, F.~A. Eckel, D.~L. Rife, B.~Nagarajan, and K.~Searight.
\newblock Probabilistic weather prediction with an analog ensemble.
\newblock \emph{Monthly Weather Review}, 141(10):3498--3516, 2013.

\bibitem{zhao2016analog}
Z.~Zhao and D.~Giannakis.
\newblock Analog forecasting with dynamics-adapted kernels.
\newblock \emph{Nonlinearity}, 29(9):2888--2939, 2016.

\bibitem{gneiting2007strictly}
T.~Gneiting and A.~E. Raftery.
\newblock Strictly proper scoring rules, prediction, and estimation.
\newblock \emph{Journal of the American Statistical Association},
 102(477):359--378, 2007.

\end{thebibliography}

\end{document}
